\scanpage{120}
\mspage{126}
\uline{Cor. 2} Supposons que $\unsure{\mathbb D|D}$ satisfasse Der 3 $(c_1,c_2)$ (\unsure{par|p.} ex. $\mathbb D=\mathrm{HOT}$). Alors \deleted{dans $\mathbb D$} pour une flèche $f:X\unsure{\mathord{-}|\to}Y$ dans Diag, il revient au $=$ que $f$ soit une \deleted{fib} $\mathbb D$-fibration, une $\mathbb D$-fibration cohomologique, ou une $\mathbb D$-fibration homologique. Si $f$ est propre (resp. lisse), \unsure{il} revient au $=$ de dire que $f$ est une $\mathbb D$-fibration coh. (resp hom.) faible.

\figure{\begin{tikzpicture}\draw (0,0) -- (.2,-.23);\end{tikzpicture}}

\deleted{\unsure{Je}}

On voit donc que dans Cat, on doit avoir une seule notion \unsure{vraisemblable|raisonnable} de HOT-fibration, qui \unsure{coïncidera} avec celle de HOT-fibration cohomologique ou homologique. Mais pour que cette notion soit dès à présent disponible, il faudrait que le dérivateur HOT soit déjà construit, sur le domaine $\mathrm{Diag}=\mathrm{Cat}$ tout entier. Or ce n'est pas chose faite encore (et il n'a \deleted{\unsure{donc}} pas pu être question de prouver les propriétés de HOT dont il vient d'être question, disons la plus forte Der 3 $c''$).

Néanmoins, je voudrais développer de façon directe (``élémentaire'') la notion de Hot-fibration, sans me référer à la théorie du dérivateur HOT, et aux développements théoriques précédents,
% archive sheet number: 120
